\section{本章小结}
\label{sec:odenum-summary}

本章回答了开头的三个问题。怎样走一步：所有方法都在近似积分式~\eqref{eq:odenum-integral-step}，显式Euler方法只看起点，Runge--Kutta方法在一步内多次采样并匹配Taylor展开，隐式方法把未知的终点也放进来。

误差怎样累积：局部截断误差沿真解定义，没有除以步长。定理~\ref{thm:odenum-global-error}说明，在解存在的有限区间上，只要方法一步稳定、一步误差一致地是$O(h^{p+1})$、初值误差和计算误差足够小，全局误差就是$O(h^{p})$。它的常数含有$\mathrm{e}^{\Lambda T}$，这是有限时间的结论，不描述长期行为。

步长由什么决定：非刚性问题由精度决定，刚性问题由稳定性决定。绝对稳定域回答固定步长下扰动是否放大，它基于线性测试方程，不能代替一般的收敛论证。

表~\ref{tab:odenum-choice}按问题的特点比较代表方法。

\begin{BookInlineTable}
 \centering
 \caption{方法选择的对比。}
 \label{tab:odenum-choice}
 \begin{tabular}{p{28mm}p{30mm}p{36mm}}
  \toprule
  问题特点 & 适用方法 & 主要代价或风险\\
  \midrule
  光滑，非刚性，低精度 & 低阶显式方法 & 精度要求提高时成本急增\\
  光滑，非刚性，高精度 & RK4或高阶嵌入对 & 需要足够的光滑性\\
  右端有已知跳变 & 在跳变点分段积分 & 不分段时阶数丢失\\
  刚性 & 后向Euler等隐式方法 & 每步解方程，内部容差计入预算\\
  解的尺度变化大 & 自适应步长 & 容差不等于全局精度\\
  \bottomrule
 \end{tabular}
\end{BookInlineTable}

两种常用的检验要分清。步长细化比较不同步长的结果，能观察实际的收敛阶，发现阶数丢失和舍入误差主导，但要求解在相应步长上已进入渐近范围。局部误差估计只度量一步新增的偏差，不包含累积，也察觉不到系统性的模型错误。

偏微分方程在空间离散后成为大型常微分方程组，时间推进可以直接使用本章的方法，空间网格越细，系统通常越刚性，具体分析见第~\ref{chap:pde-numerical-methods}章。长时间积分时的结构保持、吸引子等离散动力学问题超出了有限区间收敛定理的范围，可参考\BookChapterRef{chap:dynamical-systems-long-term}{卷九第1章“动力系统与轨迹的长期行为”}。
